\chapter{Building a nutrient-dense eating pattern}
\section{Learning objectives}
Translate nutrient science into culturally appropriate meals; identify barriers to adequacy; plan food substitutions; and avoid reducing health to a single nutrient or ``superfood.''

\section{A flexible plate}
Across the day, include vegetables or fruit, protein-rich foods, whole grains or other starchy staples, and sources of unsaturated fats in combinations that fit your culture, appetite, budget, and health needs. Rotate foods rather than searching for one perfect ingredient. Beans, lentils, tofu, eggs, fish, dairy or appropriately fortified alternatives, nuts, seeds, vegetables, fruit, and whole grains provide a broad nutrient base. \cite{who}

\section{Practical upgrades}
Diet quality is a pattern rather than a checklist. Adequacy depends on total energy, protein, fiber, essential fatty acids, and micronutrients. Health is also influenced by sodium, added sugars, alcohol, and physical activity. Processing alone does not determine whether a food fits a healthy pattern; for example, frozen vegetables and low-sodium canned beans can be practical choices. \cite{who}
\begin{itemize}
\item Add beans or lentils to soups, salads, and grain dishes for folate, iron, magnesium, potassium, and protein.
\item If your household uses salt, choose iodized salt in amounts consistent with local sodium guidance; do not add salt just to obtain iodine. Specialty salts are often not iodized. \cite{whoSaltIodization}
\item Pair plant iron sources with peppers, citrus, tomatoes, or other vitamin-C foods.
\item Choose fortified foods deliberately, especially for B12, vitamin D, calcium, and iodine when your diet excludes common sources. Check labels because products such as plant-based beverages vary in which nutrients, if any, they provide.
\item Preserve variety across the week: dark greens, orange produce, legumes, whole grains, nuts and seeds, and protein foods.
\end{itemize}

\section{Food safety and supplements}
Cooking and storage can change vitamin content, but raw is not always better. Rinse produce under running water rather than using soap, keep it separate from raw meat and seafood, and cook animal foods to safe temperatures. During pregnancy, fish can be part of a healthy diet; choose lower-mercury species and follow local fish-safety advice. Supplements may help address specific nutrient gaps, but they do not replace adequate food, sleep, or needed medical care. \cite{fdaProduce,fdaSafeFood,fdaFishPregnancy}

\section{A one-day template}
Breakfast might include fortified oatmeal, fruit, and nuts; lunch could be a lentil and vegetable bowl with yogurt or a fortified alternative; dinner could combine fish, tofu, or beans with vegetables and a whole grain. Snacks can include fruit, roasted chickpeas, or seeds. This is a template, not a prescription: adapt it to culture, budget, allergies, ethics, and medical needs.
